\clearpage
\section{Ra quyết định và kết luận}
\subsection{Tư vấn chiến lược kinh doanh}
%Dành cho Người Bán / Nhà đầu tư (House Flippers): Tập trung nâng cấp yếu tố nào để tối ưu ROI (ví dụ: chất lượng vật liệu thay vì cơi nới diện tích).
%Dành cho Người Mua nhà: Lời khuyên tối ưu ngân sách khi chọn mua nhà.
Từ những tri thức được khai phá từ mô hình học máy, nhóm đề xuất các chiến lược kinh doanh thực tiễn cho các đối tượng tham gia thị trường bất động sản:

\begin{itemize}
    \item[-] \textbf{Dành cho Người Bán và Nhà đầu tư (House Flippers):} Thuật toán chỉ ra rằng \texttt{OverallQual} (Chất lượng vật liệu và hoàn thiện) tác động đến giá bán mạnh hơn diện tích. Thay vì tốn chi phí và thời gian xin giấy phép để cơi nới không gian sinh hoạt (\texttt{GrLivArea}), nhà đầu tư nên tối ưu tỷ suất hoàn vốn (ROI) bằng cách phân bổ ngân sách cho việc nâng cấp ngoại thất, tân trang nội thất, và sử dụng vật liệu cao cấp. Một căn nhà diện tích vừa phải nhưng có thiết kế và chất lượng nội thất xuất sắc sẽ mang lại biên độ lợi nhuận cao hơn.
    
    \item[-] \textbf{Dành cho Người Mua nhà:} Đối với những gia đình hoặc cá nhân có ngân sách hạn hẹp nhưng muốn tối ưu hóa không gian sống, chiến lược tốt nhất là tìm kiếm những căn nhà có diện tích lớn (\texttt{GrLivArea} cao) nhưng chất lượng hiện tại ở mức trung bình hoặc thấp (\texttt{OverallQual} thấp). Vì những căn nhà này bị định giá thấp hơn rất nhiều so với giá trị không gian thực của chúng, người mua có thể chốt giao dịch với giá hời, sau đó tự thực hiện quá trình cải tạo và tân trang dần dần để nâng tầm giá trị căn nhà trong tương lai.
\end{itemize}

\subsection{Kết luận}
Dự án đã hoàn thành các mục tiêu được đặt ra ở Chương 1. Nhóm đã xây dựng thành công một pipelive khai phá dữ liệu hoàn chỉnh: từ khâu tiền xử lý, làm sạch và loại bỏ nhiễu, đến việc huấn luyện mô hình cơ sở (Linear Regression) và mô hình nâng cao (Random Forest). \medskip

Mô hình Random Forest đã chứng minh được sự vượt trội trong việc xử lý tập dữ liệu nhà ở nhiều chiều, giảm thiểu đáng kể sai số dự đoán so với mô hình cơ sở. Quan trọng hơn, thông qua kỹ thuật giải thích mô hình, đồ án đã chuyển hóa các con số khô khan thành những tư vấn chiến lược mang tính thực tiễn cao, hỗ trợ trực tiếp quá trình ra quyết định cho cả người mua và người bán.

\subsection{Hướng phát triển tương lai}
Dù mô hình hiện tại đã đạt được độ chính xác tốt, hệ thống vẫn còn nhiều không gian để cải thiện và mở rộng nếu có thêm thời gian và nguồn lực:

\begin{itemize}
    \item[-] \textbf{Thử nghiệm các thuật toán tiên tiến hơn:} Triển khai các thuật toán Gradient Boosting mạnh mẽ như XGBoost, LightGBM hoặc CatBoost, kết hợp với các kỹ thuật tinh chỉnh siêu tham số (Hyperparameter Tuning bằng GridSearchCV hoặc Optuna) để tiếp tục ép mức sai số dự đoán xuống mức thấp nhất.
    
    \item[-] \textbf{Bổ sung dữ liệu không gian và kinh tế vĩ mô:} Giá bất động sản chịu ảnh hưởng lớn bởi yếu tố ngoại cảnh. Hệ thống có thể được làm giàu dữ liệu bằng cách thu thập thêm các thông tin về tọa độ địa lý (khoảng cách đến trung tâm, trường học, bệnh viện) hoặc các chỉ số kinh tế (lãi suất vay thế chấp tại thời điểm bán).
    
    \item[-] \textbf{Phát triển ứng dụng Web (End-to-End):} Tích hợp mô hình đã huấn luyện (pickle/joblib) vào một hệ thống Web API (sử dụng FastAPI, Flask, hoặc tích hợp Node.js/React theo hướng quan tâm của nhóm) để người dùng có thể trực tiếp nhập thông số và nhận định giá theo thời gian thực.
\end{itemize}
